\section{Definición de ciclo celular}

El ciclo celular es el ciclo vital de una célula y es una serie de etapas de desarrollo y crecimiento entre su nacimiento y reproducción.

\section{Fases del ciclo celular}

\begin{enumerate}
    \item \textbf{(M) Fase mitótica}: Se separa y divide el ADN para formar dos células
    \item \textbf{(I) Interfase}: Crecimiento celular
    \item \textbf{(G)}:
        \begin{itemize}
            \item G0: Control - La célula se encuentra en reposo
            \item G1: Crecimiento físico y preparación para la replicación
            \item S: Replicación del ADN (Síntesis)
            \item G2: Más crecimiento y síntesis de proteínas y organelos
        \end{itemize}
\end{enumerate}

\section{Tipos de división celular}

\subsection{Mitosis}
La mitosis es la división de células somáticas (células del cuerpo no reproductivas).

\begin{itemize}
    \item \textbf{Células somáticas}: Todas las células del cuerpo excepto gametos
    \item \textbf{Característica}: Las células "hijas" son idénticas a la célula "madre" (2n → 2n)
    \item \textbf{Fases}:
        \begin{itemize}
            \item \textbf{Profase}: Condensación de filamentos de cromatina para generar cromosomas. Desaparece el nucléolo y aparece el huso mitótico.
            \item \textbf{Metafase}: Alineación de los cromosomas en el plano ecuatorial por los centrosomas.
            \item \textbf{Anafase}: Separación de las cromátidas hermanas hacia polos opuestos.
            \item \textbf{Telofase}: Descondensación de cromosomas, reorganización del núcleo y citocinesis ("cierre" de las células)
        \end{itemize}
\end{itemize}

\subsection{Meiosis}
La meiosis es la división que produce gametos (células reproductivas).

\begin{itemize}
    \item \textbf{Gametos}: Células sexuales (óvulos y espermatozoides)
    \item \textbf{Característica}: Dos divisiones celulares consecutivas que reducen el material genético a la mitad (2n → n)
    \item \textbf{Importancia}: Variabilidad genética y formación de gametos haploides
\end{itemize}

\section{Comparación entre mitosis y meiosis}

\begin{table}[h]
\centering
\begin{tabular}{|l|c|c|}
\hline
\textbf{Característica} & \textbf{Mitosis} & \textbf{Meiosis} \\
\hline
Tipo de célula & Somática & Gametos \\
\hline
Número de divisiones & 1 & 2 \\
\hline
Células hijas & 2 & 4 \\
\hline
Ploidía & 2n → 2n & 2n → n \\
\hline
Identidad genética & Idénticas & Diferentes \\
\hline
Crossing-over & No & Sí \\
\hline
\end{tabular}
\caption{Comparación entre mitosis y meiosis}
\end{table}

\section{Regulación del ciclo celular}

\begin{itemize}
    \item \textbf{Puntos de control}: G1/S, G2/M, M
    \item \textbf{Quinasas dependientes de ciclina (CDK)}: Reguladores clave
    \item \textbf{Oncogenes y genes supresores de tumores}: Importantes en el control
\end{itemize}